% Generated draft: requires mathematical and source review.
\documentclass[11pt]{article}
\usepackage[margin=25mm]{geometry}
\usepackage{fontspec}
\setmainfont{Latin Modern Roman}
\usepackage{amsmath,amssymb,mathtools,unicode-math}
\setmathfont{Latin Modern Math}
\usepackage{enumitem,needspace,xurl,hyperref}
\setlength{\parindent}{0pt}
\setlength{\parskip}{5pt}
\setlength{\emergencystretch}{3em}
\newcommand{\N}{\mathbb{N}}
\newcommand{\Z}{\mathbb{Z}}
\newcommand{\Q}{\mathbb{Q}}
\newcommand{\R}{\mathbb{R}}
\newcommand{\C}{\mathbb{C}}
\newcommand{\F}{\mathbb{F}}
\newcommand{\A}{\mathbb{A}}
\newcommand{\PP}{\mathbb{P}}
\newcommand{\E}{\mathbb{E}}
\newcommand{\eps}{\varepsilon}
\newcommand{\dd}{\,\mathrm{d}}
\begin{document}
\section*{UnsolvedMath ID: 30002287}
\textbf{Source problem code:} catalogue entry
\subsection{Definitions and mathematical statement}
For a conventional formulation, let \(n\geq 1\), let \(\Omega\subset\R^n\) be a nonempty bounded connected open set, and let \(1<p<\infty\). These parameter ranges and the connectedness convention are not specified explicitly in the supplied record; any additional intended hypotheses on \(\Omega\) require verification. Write \(W^{1,p}_0(\Omega)\) for the closure of \(C_c^\infty(\Omega)\) in \(W^{1,p}(\Omega)\). The Dirichlet \(p\)-Laplacian eigenvalue equation is
\[
-\operatorname{div}(|\nabla u|^{p-2}\nabla u)=\lambda |u|^{p-2}u,
\qquad u\in W^{1,p}_0(\Omega)\setminus\{0\}.
\]
The equation is understood weakly: for every \(\varphi\in W^{1,p}_0(\Omega)\),
\[
\int_\Omega |\nabla u|^{p-2}\nabla u\cdot\nabla\varphi\,\dd x
=\lambda\int_\Omega |u|^{p-2}u\varphi\,\dd x.
\]
Define the full eigenvalue set by
\[
\Sigma_p(\Omega)=\{\lambda\in\R:\text{there exists a nonzero weak solution of the displayed equation}\}.
\]
This definition includes all weak eigenvalues, without assuming that a variational sequence exhausts them.

\textbf{Discreteness question.} In the conventional framework above, determine whether, for every admissible \(\Omega\) and \(p\), the full set \(\Sigma_p(\Omega)\) has no finite accumulation point. An explicit formulation of this interpretation is
\[
\forall M>0,\qquad \#\bigl(\Sigma_p(\Omega)\cap[0,M]\bigr)<\infty.
\]
The record does not define ``discrete.'' If only isolation of each eigenvalue is intended, the corresponding assertion is instead
\[
\forall\lambda\in\Sigma_p(\Omega)\ \exists\delta>0:\quad
\Sigma_p(\Omega)\cap(\lambda-\delta,\lambda+\delta)=\{\lambda\}.
\]
The intended convention should be checked against the report.

\textbf{Further questions recorded by the source.} Determine whether higher eigenvalues have finite multiplicity, and whether there exists
\[
\mu>\lambda_2\quad\text{such that}\quad\mu\notin\Sigma_p(\Omega).
\]
For an explicit conventional meaning of \(\lambda_2\), put
\[
S=\left\{u\in W^{1,p}_0(\Omega):\int_\Omega |u|^p\,\dd x=1\right\},
\qquad E(u)=\int_\Omega |\nabla u|^p\,\dd x,
\]
choose a positive minimizer \(u_1\in S\) of \(E\), and set
\[
\lambda_2=\inf_{\gamma\in\Gamma}\ \max_{t\in[0,1]}E(\gamma(t)),\qquad
\Gamma=\{\gamma\in C([0,1],S):\gamma(0)=u_1,\ \gamma(1)=-u_1\},
\]
where continuity uses the \(W^{1,p}_0(\Omega)\) topology. This supplies a conventional second variational eigenvalue; the supplied record does not state its convention explicitly. The meaning of ``multiplicity'' remains unresolved: for \(p\neq2\), the solutions associated with a fixed eigenvalue need not form a vector space, so ordinary eigenspace dimension cannot be adopted without justification. Accordingly, the finite-multiplicity question is retained as a source question requiring a specified nonlinear multiplicity notion, rather than asserted to follow from discreteness of the eigenvalue set.

The standard equation above uses the same exponent \(p\) in the operator and in \(|u|^{p-2}u\). The supplied assessment distinguishes it from a \((p,q)\)-problem with \(q<p\); the reported non-discrete example for that different problem is not presented as an answer to this question.
\subsection{Short English statement}
The problem asks whether the full Dirichlet \(p\)-Laplacian eigenvalue set on a bounded domain is discrete. It also asks whether higher eigenvalues have finite multiplicity and whether at least one real number above the second eigenvalue fails to be an eigenvalue. The precise domain hypotheses, meaning of discreteness, and nonlinear multiplicity convention need confirmation. The supplied registry labels the problem open; this is dated curation, not an independently verified status determination.
\subsection{Sources}
\begin{itemize}
\item Local UnsolvedMath registry, permanent ID \(30002287\), source problem code \texttt{catalogue entry}, title: \emph{Discreteness of the Dirichlet \(p\)-Laplacian Spectrum}. The record belongs to the Oberwolfach Reports Open Problems collection and gives proposed year 2013. No individual proposer is supplied. Registry locations identified by the project instructions are \texttt{data/unsolvedmath/problems.json} and the supplementary \texttt{catalog/problems.json}; neither was independently inspected for this draft.
\item Supplied primary attribution: \emph{Mini-Workshop: The p-Laplacian Operator and Applications} (2013), Oberwolfach Reports, pp.~433--482, DOI \texttt{10.4171/owr/2013/08}. \url{https://doi.org/10.4171/owr/2013/08}. The supplied background describes these reports as records of mathematical workshops and their open-problem discussions.
\item Supplied literature source: Le, \emph{Eigenvalue problems for the p-Laplacian}. \url{https://doi.org/10.1016/j.na.2005.05.056}.
\item Supplied literature source: Brasco--Franzina, described in the registry as concerning a non-discrete \((p,q)\)-spectrum. \url{https://arxiv.org/abs/1707.00867}.
\item Literature assessment date: \texttt{2026-08-21}, copied exactly from \texttt{literature\_checked\_at}. This is supplied dated curation, not independently verified. It labels the problem open and states that discreteness of the standard Dirichlet \(p\)-Laplacian spectrum on general bounded domains remains unknown, while the known non-discrete example concerns the distinct \((p,q)\)-problem with \(q<p\). The background explicitly requests independent verification of time-sensitive status and cited claims.
\item The supplied statement status is \texttt{corrected\_verified}. Its verification note says that the official report identifies discreteness as the main open problem and records both additional questions as weaker unknown consequences. This is a supplied source-verification claim; the report and linked literature were not independently consulted for this draft.
\end{itemize}
\end{document}
